\chapter{Dette en difficulté, dette souveraine et crédit en fonds propres bancaires}\label{ch:m2:distressed-sovereign-and-bank-capital-credit}

Le dimanche 19 mars 2023, les autorités suisses ont annoncé qu'UBS reprendrait Credit
Suisse. Le même jour, le superviseur, la FINMA, a indiqué que le soutien public exceptionnel
requis par l'opération avait entraîné une dépréciation totale de toutes les obligations
\omterm{def:m2:distressed-sovereign-and-bank-capital-credit:coco}{Additional Tier 1} (AT1) de Credit Suisse, soit environ 16~milliards de CHF, ramenées à zéro. Les actionnaires
de Credit Suisse, qui passent après ces obligations en cas de liquidation, ont reçu des actions UBS valant environ 3~milliards de CHF.
Le lendemain, les superviseurs et l'autorité de résolution européens ont affirmé que dans
l'Union européenne, les fonds propres absorbent les pertes en premier. En octobre 2025, un tribunal suisse a jugé
la dépréciation illégale~; la FINMA a fait appel. Ce chapitre traite du crédit une fois que les choses ont mal tourné~:
la \omterm{def:m2:distressed-sovereign-and-bank-capital-credit:distressed}{dette en difficulté} et les \omterm{def:m2:distressed-sovereign-and-bank-capital-credit:distressed}{offres d'échange}, les souverains qu'on ne peut pas traduire
devant un tribunal de faillite, et les obligations bancaires conçues pour absorber les pertes avant la défaillance de la banque.

\section{Restructuration et offres d'échange}

\begin{definition}[Dette en difficulté, offre d'échange]\label{def:m2:distressed-sovereign-and-bank-capital-credit:distressed}
Une \emph{dette en difficulté}\index{dette en difficulté} (distressed debt) est une dette d'un emprunteur en défaut ou proche du défaut,
cotée à un prix qui traduit une perte espérée plutôt qu'un spread au-dessus du
taux sans risque. Une \emph{offre d'échange}\index{offre d'échange} (exchange offer) est une proposition de l'emprunteur à ses créanciers
d'échanger leurs créances contre de nouveaux instruments, en général de moindre valeur, avec
une valeur nominale plus faible, des coupons plus bas ou des échéances plus longues, parfois avec un versement en espèces.
\end{definition}

Une entreprise en difficulté peut se placer en faillite, procédure dans laquelle un tribunal impose un plan à tous
les créanciers par catégorie~; ou bien elle négocie à l'amiable, et c'est par des \omterm{def:m2:distressed-sovereign-and-bank-capital-credit:distressed}{offres d'échange} qu'elle le fait.
À l'amiable, la voie peut être moins coûteuse et plus rapide, mais chaque créancier décide pour lui-même, et un
créancier qui n'accepte pas conserve son ancienne créance. La valeur d'une offre est la valeur actuelle
de ce qu'elle procure, au rendement auquel les nouveaux instruments se négocieront une fois
l'échange, le \emph{rendement de sortie} (exit yield), comparée à l'ancienne créance~: la perte
en valeur actuelle est la décote de l'offre, souvent bien supérieure à la baisse de la valeur nominale.

\begin{example}[Une offre d'échange]\label{ex:m2:distressed-sovereign-and-bank-capital-credit:offer}
Un emprunteur propose, pour 100 d'anciennes obligations, 50 de nouvelles obligations de 15 ans rapportant 4~\% et 5 en
espèces pour ceux qui apportent leurs titres. Avec un rendement de sortie de 9~\%, les nouvelles obligations valent 29.85 et
le lot 34.85~: une décote en valeur nominale de 50~\% mais une décote en valeur actuelle de 65.2~\%.
Avec un rendement de sortie de 6~\%, le lot vaut 45.29, à 12~\% seulement 27.76
(voir la \cref{fig:m2:distressed-sovereign-and-bank-capital-credit:offer}).
\end{example}

\begin{omfigure}
\begin{tikzpicture}
\begin{axis}[width=11cm,height=4.8cm,
  xlabel={rendement de sortie (\%)},ylabel={valeur pour 100 de nominal ancien},
  tick label style={font=\small},label style={font=\small},
  xmin=5,xmax=15,ymin=15,ymax=60,grid=major,grid style={gray!25},
  legend columns=2,legend style={font=\footnotesize,at={(0.5,-0.3)},anchor=north,draw=none,fill=none,/tikz/every even column/.append style={column sep=8pt}},
  table/col sep=comma]
\addplot[omDef,very thick] table[x=yield,y=withcash]{figdata/markets-2/26-distressed-sovereign-and-bank-capital-credit/offer.csv};
\addplot[omThm,very thick,dashed] table[x=yield,y=bonds]{figdata/markets-2/26-distressed-sovereign-and-bank-capital-credit/offer.csv};
\addplot[only marks,mark=*,mark size=1.6pt,black,forget plot] coordinates {(6,45.29) (9,34.85) (12,27.76)};
\legend{nouvelles obligations et espèces,nouvelles obligations seules}
\end{axis}
\end{tikzpicture}

\omcaption{Valeur de l'\omterm{def:m2:distressed-sovereign-and-bank-capital-credit:distressed}{offre d'échange} du chapitre, pour 100 de nominal ancien, en fonction du
rendement de sortie auquel les nouvelles obligations de 15 ans à 4~\% se négocieront~: avec les 5 en espèces versés à ceux
qui apportent leurs titres, et sans ce versement, ce qui correspond à ce que reçoit un détenteur lié par une \omterm{def:m2:distressed-sovereign-and-bank-capital-credit:cac}{clause d'action collective}.
reçoit. Les points représentent l'\cref{ex:m2:distressed-sovereign-and-bank-capital-credit:offer}.
Illustratif~; données~: tutoriel du chapitre.}\label{fig:m2:distressed-sovereign-and-bank-capital-credit:offer}
\end{omfigure}

\section{Dette souveraine~: clauses d'action collective et créanciers récalcitrants}

Un État souverain n'a pas de tribunal de faillite. Ses obligations ne peuvent être restructurées que par accord,
et chaque créancier qui n'accepte pas peut refuser, conserver sa créance et engager des poursuites.

\begin{definition}[Clause d'action collective, créancier récalcitrant]\label{def:m2:distressed-sovereign-and-bank-capital-credit:cac}
Une \emph{clause d'action collective}\index{clause d'action collective} (collective action clause, CAC) dans une obligation permet à une
majorité qualifiée de ses détenteurs de modifier ses conditions de paiement pour tous les détenteurs, y compris
la minorité. Un \emph{créancier récalcitrant}\index{créancier récalcitrant} (holdout creditor) est un créancier qui
n'accepte pas une restructuration et conserve, ou poursuit en justice, sa créance d'origine.
\end{definition}

Le FMI décrit la clause traditionnelle comme s'appliquant série par série~: en général, 75~\% des
détenteurs d'une même obligation peuvent la modifier, de sorte qu'un créancier qui achète une minorité de blocage dans une série
peut tenir cette série à l'écart. Les clauses plus récentes agrègent les séries~: dans les clauses types
de l'International Capital Market Association, un vote unique de 75~\% du principal de toutes les
séries concernées engage toutes les séries, à condition que toutes soient offertes aux mêmes conditions, ou,
dans une version à deux volets, 66\,2/3~\% au total et plus de la moitié de chaque série.

\begin{definition}[Clause pari passu]\label{def:m2:distressed-sovereign-and-bank-capital-credit:pp}
Une \emph{clause pari passu}\index{clause pari passu} (pari passu clause) stipule qu'une obligation a un rang égal
avec les autres dettes chirographaires de l'émetteur de même nature.
\end{definition}

Deux cas ont façonné la pratique actuelle. L'Argentine a fait défaut en 2001 et a restructuré environ
93~\% de sa dette extérieure par des échanges en 2005 et 2010. Des \omterm{def:m2:distressed-sovereign-and-bank-capital-credit:cac}{créanciers récalcitrants} menés par NML Capital,
pour des créances d'environ 1.6~milliard de dollars, ont engagé des poursuites à New York et ont obtenu une interprétation de la \omterm{def:m2:distressed-sovereign-and-bank-capital-credit:pp}{clause pari
passu} comme exigeant un paiement au prorata~: l'Argentine ne pouvait pas payer les obligations échangées
sans payer intégralement dans le même temps les \omterm{def:m2:distressed-sovereign-and-bank-capital-credit:cac}{créanciers récalcitrants}, et les banques de la chaîne de paiement
se voyaient interdire de transmettre ces paiements. La Cour suprême a refusé d'examiner l'appel de l'Argentine
en juin 2014, et fin juillet 2014 les détenteurs des obligations échangées, qui avaient
accepté environ 27 cents pour un dollar en valeur actuelle, n'ont pas été payés. La Grèce en 2012
a fait le contraire~: une loi adoptée en février a introduit des \omterm{def:m2:distressed-sovereign-and-bank-capital-credit:cac}{clauses d'action collective} dans
ses obligations de droit grec, qui les rendaient contraignantes si les deux tiers des détenteurs participants acceptaient. Les détenteurs
d'environ 85~\% des 206~milliards d'euros d'obligations détenues par le secteur privé ont accepté, les clauses
ont rendu l'échange contraignant pour les obligations de droit grec, et de nouvelles obligations d'une valeur nominale d'environ 100~milliards d'euros
de valeur nominale ont remplacé les anciennes. Environ 6~milliards d'euros d'obligations de droit anglais sont restées hors de l'échange~; en
mai 2012, la Grèce a remboursé intégralement une obligation de droit anglais de 435~millions d'euros arrivant à échéance.

Pourquoi rester à l'écart~? La valeur d'une position de \omterm{def:m2:distressed-sovereign-and-bank-capital-credit:cac}{créancier récalcitrant} est un arbre de probabilités
(voir la \cref{fig:m2:distressed-sovereign-and-bank-capital-credit:tree})~: l'État peut
rembourser intégralement les \omterm{def:m2:distressed-sovereign-and-bank-capital-credit:cac}{créanciers récalcitrants} pour s'en débarrasser, ce qui est d'autant plus probable qu'ils sont peu
nombreux, comme l'a fait la Grèce~; sinon, une action en justice peut obtenir des années plus tard la totalité de la créance, comme NML l'a obtenu~; ou
l'obligation peut rester en défaut. Plus les autres acceptent l'échange, plus la position de \omterm{def:m2:distressed-sovereign-and-bank-capital-credit:cac}{créancier récalcitrant} paraît attrayante,
ce qui est le problème du passager clandestin que les CAC ont pour but de résoudre.

\begin{omfigure}
\centering
\begin{tikzpicture}[font=\small,
  nd/.style={draw,thick,rounded corners=2pt,align=center,inner sep=3pt,minimum height=0.8cm},
  arr/.style={-{Stealth[length=2.2mm]},thick}]
\node[nd,draw=omDef,fill=omDef!8] (h) at (0,0) {rester à l'écart\\(participation $p$)};
\node[nd,draw=omThm,fill=omThm!8] (paid) at (5,1.4) {remboursé intégralement l'année suivante\\valeur 91.74};
\node[nd,draw=gray,fill=gray!8] (no) at (5,-1.2) {non remboursé à court terme};
\node[nd,draw=omProp,fill=omProp!8] (law) at (10.4,-0.3) {l'action en justice récupère 130\\en six ans~: 77.51};
\node[nd,draw=gray,fill=gray!15] (stuck) at (10.4,-2.1) {reste en défaut\\valeur 15};
\draw[arr] (h) -- node[above,sloped,font=\footnotesize] {$p^8$} (paid);
\draw[arr] (h) -- node[below,sloped,font=\footnotesize] {$1 - p^8$} (no);
\draw[arr] (no) -- node[above,sloped,font=\footnotesize] {0.3} (law);
\draw[arr] (no) -- node[below,sloped,font=\footnotesize] {0.7} (stuck);
\end{tikzpicture}

\omcaption{Arbre de probabilités du \omterm{def:m2:distressed-sovereign-and-bank-capital-credit:cac}{créancier récalcitrant} dans l'exemple du chapitre, valeurs pour 100 de l'ancienne
valeur faciale actualisée à 9~\%. La probabilité d'être remboursé intégralement à court terme augmente fortement avec la part
des créanciers ayant apporté leurs titres, puisque le souverain peut se permettre de payer un petit
reliquat. Illustratif.}\label{fig:m2:distressed-sovereign-and-bank-capital-credit:tree}
\end{omfigure}

\begin{proposition}[La bande des passagers clandestins]\label{prop:m2:distressed-sovereign-and-bank-capital-credit:band}
Soit $T$ la valeur de l'apport des titres, $H(p)$ la valeur de la rétention lorsqu'une part $p$
des créanciers apporte ses titres, croissante en $p$, et $p^*$ la participation pour laquelle $H(p^*) = T$. Sans \omterm{def:m2:distressed-sovereign-and-bank-capital-credit:cac}{clause
d'action collective}, la rétention rapporte dès que $p > p^*$. Avec une clause agrégée
de seuil $\tau > p^*$, un \omterm{def:m2:distressed-sovereign-and-bank-capital-credit:cac}{créancier récalcitrant} est lié par l'offre dès que $p \geq \tau$ et ne reçoit
que ce que l'offre accorde au créancier lié~; la rétention ne rapporte que pour $p^* < p < \tau$, et pas
du tout si $\tau \leq p^*$.
\end{proposition}

\begin{proof}
En dessous de $\tau$, le \omterm{def:m2:distressed-sovereign-and-bank-capital-credit:cac}{créancier récalcitrant} conserve sa créance, d'une valeur $H(p)$, qui dépasse $T$ exactement lorsque
$p > p^*$. À partir de $\tau$ et au-delà, la clause modifie son titre pour l'aligner sur les conditions de l'offre, lesquelles
valent au plus $T$ et moins toute somme en numéraire réservée à ceux qui ont apporté leurs titres.
\end{proof}

\begin{example}[Où la rétention rapporte]\label{ex:m2:distressed-sovereign-and-bank-capital-credit:band}
Dans l'arbre de la \cref{fig:m2:distressed-sovereign-and-bank-capital-credit:tree}, la rétention
vaut 33.75 si personne n'apporte ses titres, 34.73 pour une participation de 60~\% et 38.97 pour 74~\%,
contre 34.85 pour l'apport. Elle rapporte à partir d'une participation de 60.9~\%. Avec une clause agrégée de 75~\%
à 75~\% ou au-delà, le \omterm{def:m2:distressed-sovereign-and-bank-capital-credit:cac}{créancier récalcitrant} est lié et ne reçoit que les nouvelles obligations,
soit 29.85~: la bande dans laquelle la rétention rapporte va de 60.9~\% à 75~\%. Sans clause,
elle rapporte pour toute participation supérieure à 60.9~\%, et à 90~\% elle vaut 58.72
(la \cref{fig:m2:distressed-sovereign-and-bank-capital-credit:holdout}).
\end{example}

\begin{omfigure}
\begin{tikzpicture}
\begin{axis}[width=11cm,height=5cm,
  xlabel={participation à l'échange (\%)},ylabel={valeur pour 100 de valeur faciale ancienne},
  tick label style={font=\small},label style={font=\small},
  xmin=0,xmax=100,ymin=20,ymax=95,grid=major,grid style={gray!25},
  legend columns=3,legend style={font=\footnotesize,at={(0.5,-0.3)},anchor=north,draw=none,fill=none,/tikz/every even column/.append style={column sep=8pt}},
  table/col sep=comma]
\addplot[omThm,very thick,dashed] table[x=p,y=nocac]{figdata/markets-2/26-distressed-sovereign-and-bank-capital-credit/holdout.csv};
\addplot[omDef,very thick,const plot] table[x=p,y=cac]{figdata/markets-2/26-distressed-sovereign-and-bank-capital-credit/holdout.csv};
\addplot[black,thick] table[x=p,y=tender]{figdata/markets-2/26-distressed-sovereign-and-bank-capital-credit/holdout.csv};
\addplot[gray,thin,dotted,forget plot] coordinates {(60.9,20) (60.9,95)};
\addplot[gray,thin,dotted,forget plot] coordinates {(75,20) (75,95)};
\node[font=\footnotesize,anchor=north east] at (axis cs:60.4,93) {$p^* = 60.9\%$};
\node[font=\footnotesize,anchor=north west] at (axis cs:75.5,93) {CAC 75\%};
\legend{rétention~: sans CAC,rétention~: CAC 75\%,apport}
\end{axis}
\end{tikzpicture}

\omcaption{Valeur de la rétention en fonction de la part des créanciers qui apportent leurs titres, sans
\omterm{def:m2:distressed-sovereign-and-bank-capital-credit:cac}{clause d'action collective} et avec une clause agrégée de 75~\%, ainsi que la valeur de l'apport. La rétention
rapporte au-delà d'une participation de 60.9~\%~; la clause supprime le gain à partir de 75~\%,
où le \omterm{def:m2:distressed-sovereign-and-bank-capital-credit:cac}{créancier récalcitrant} est lié par l'offre et perd la somme en numéraire versée à ceux qui ont apporté leurs titres.
Illustratif~; données~: le tutoriel du chapitre.}\label{fig:m2:distressed-sovereign-and-bank-capital-credit:holdout}
\end{omfigure}

\section{Capital bancaire~: obligations convertibles contingentes}

Une banque qui fait défaut lors d'une crise peut être trop importante pour être liquidée. Après 2008,
la réponse des régulateurs a consisté à faire absorber des pertes à certains de ses créanciers tandis que la banque
poursuit son activité.

\begin{definition}[Obligation convertible contingente, Additional Tier 1]\label{def:m2:distressed-sovereign-and-bank-capital-credit:coco}
Une \emph{obligation convertible contingente}\index{obligation convertible contingente} (contingent \omterm{def:m2:emerging-market-fx-and-ndfs:controls}{convertible} bond, CoCo) est une
obligation bancaire qui est convertie en actions ou dépréciée lorsqu'un déclencheur est atteint, comme
le ratio de fonds propres ordinaires de la banque passe sous un niveau fixé dans l'obligation. Le capital \emph{Additional
Tier 1}\index{Additional Tier 1} (AT1) se compose d'instruments perpétuels,
généralement des CoCo, dont la banque peut annuler les coupons et qui absorbent des pertes tant que la
banque est en activité~; ils sont classés au-dessus des fonds propres ordinaires et en dessous de toute autre dette en cas de
liquidation.
\end{definition}

\begin{definition}[Point de non-viabilité, renflouement interne]\label{def:m2:distressed-sovereign-and-bank-capital-credit:pon}
Le \emph{point de non-viabilité}\index{point de non-viabilité} est le moment où
une autorité décide qu'une banque ferait défaut sans une dépréciation de ses instruments de capital
ou un soutien public. Le \emph{renflouement interne}\index{renflouement interne} (bail-in) est la dépréciation ou
la conversion en actions des passifs d'une banque défaillante par une autorité, de sorte que
les créanciers plutôt que les contribuables supportent ses pertes.
\end{definition}

Le Comité de Bâle a exigé en janvier 2011 que, à partir de 2013, tous les instruments de capital
autres que les fonds propres ordinaires émis par les banques actives à l'international puissent être dépréciés ou
convertis en fonds propres ordinaires, à la discrétion de l'autorité, au plus tôt à la date où l'on décide
qu'une dépréciation est nécessaire pour maintenir la banque viable, ou à la date où l'on décide d'injecter des fonds publics
ou un soutien équivalent sans lequel elle ne le serait pas. Une obligation AT1
comporte donc deux déclencheurs~: le sien, lié au ratio, tant que la banque est en activité, et
celui de l'autorité au \omterm{def:m2:distressed-sovereign-and-bank-capital-credit:pon}{point de non-viabilité}. Ses détenteurs touchent un coupon élevé en rémunération
du risque que l'un ou l'autre se réalise, et du risque que la banque cesse simplement de payer le
coupon, ce qu'elle peut faire sans être en défaut.

\begin{omfigure}
\centering
\begin{tikzpicture}[font=\small]
\def\w{3.2}
\foreach \x/\t in {0/{ordre du communiqué de l'UE},6.4/{Credit Suisse, 2023}} \node[font=\small\bfseries] at (\x+\w/2,5.3) {\t};
\foreach \x in {0,6.4} {
  \fill[gray!12] (\x,3.2) rectangle (\x+\w,4.8); \node[align=center] at (\x+\w/2,4.0) {dette senior\\(bail-in)};
  \fill[omProp!20] (\x,2.4) rectangle (\x+\w,3.2); \node at (\x+\w/2,2.8) {Tier 2};
  \fill[omThm!25] (\x,1.2) rectangle (\x+\w,2.4); \node at (\x+\w/2,1.8) {AT1};
  \fill[omDef!25] (\x,0) rectangle (\x+\w,1.2); \node[align=center] at (\x+\w/2,0.6) {fonds propres\\ordinaires};
  \draw[thick] (\x,0) rectangle (\x+\w,4.8);
  \foreach \y in {1.2,2.4,3.2} \draw[thick] (\x,\y) -- (\x+\w,\y);
}
\draw[-{Stealth[length=2.5mm]},very thick,omDef] (3.5,0.2) -- node[right,font=\footnotesize,align=left] {1. d'abord les\\fonds propres} (3.5,1.1);
\draw[-{Stealth[length=2.5mm]},very thick,omThm] (3.5,1.3) -- node[right,font=\footnotesize,align=left] {2. puis\\l'AT1} (3.5,2.3);
\draw[-{Stealth[length=2.5mm]},very thick,omThm] (9.9,1.3) -- node[right,font=\footnotesize,align=left] {AT1 ramenée\\à zéro} (9.9,2.3);
\node[font=\footnotesize,align=left,anchor=west] at (9.8,0.6) {actionnaires\\versés 3~Md~CHF};
\end{tikzpicture}

\omcaption{La pile d'absorption des pertes d'une banque. Les autorités européennes ont déclaré en mars 2023
que les fonds propres ordinaires absorbent les pertes en premier et l'AT1 seulement après leur épuisement complet~; dans le cas de Credit
Suisse, les obligations AT1 ont été réduites à zéro en vertu de leur déclencheur contractuel de non-viabilité
et d'une ordonnance d'urgence, tandis que les actionnaires recevaient des actions UBS. Schéma~; hauteurs
non à l'échelle.}\label{fig:m2:distressed-sovereign-and-bank-capital-credit:stack}
\end{omfigure}

\section{La dépréciation de 2023}

Les obligations AT1 de Credit Suisse prévoyaient une dépréciation intégrale
en cas d'\emph{événement de non-viabilité}, en particulier si la banque recevait un soutien public
extraordinaire. Le 19 mars 2023, elle a reçu des prêts d'aide de liquidité adossés à une garantie fédérale
contre le défaut, et une ordonnance d'urgence du Conseil fédéral a autorisé la FINMA à ordonner la réduction à zéro
des AT1~; la FINMA l'a ordonnée, et la valeur nominale des obligations AT1
d'environ 16~milliards de CHF est devenue nulle, ce qui a relevé les fonds propres de base de la banque. Les obligations de rang 2 n'ont pas
été dépréciées. Selon les termes de la fusion, les actionnaires de Credit Suisse ont reçu une action UBS
pour chaque tranche de 22.48 actions détenues, soit 3~milliards de CHF au total
(\cref{fig:m2:distressed-sovereign-and-bank-capital-credit:stack}).

Pour les créanciers obligataires qui lisaient la pile de capital comme un ordre de liquidation, il s'agissait d'un renversement~:
la créance qui prime les actions n'a rien reçu, tandis que les actions ont reçu quelque chose. La
déclaration des superviseurs européens du 20 mars, selon laquelle les fonds propres absorbent les pertes en premier et l'AT1
seulement après leur épuisement complet, a fait savoir aux investisseurs que le cadre européen ne ferait pas la
même chose. La dépréciation a suivi les termes contractuels des obligations plutôt qu'un ordre de
liquidation, et la question de savoir si ces termes et l'ordonnance la permettaient est devenue une affaire
judiciaire~: le 1\up{er} octobre 2025, le Tribunal administratif fédéral, dans la première d'environ 360
procédures, a jugé que l'événement de non-viabilité contractuel ne s'était pas produit et que la
dépréciation n'avait pas de base juridique~; le jugement n'est pas définitif, et la FINMA a fait appel devant le
Tribunal fédéral.

\begin{dated}{2026-09}{Le contentieux AT1 de Credit Suisse}\label{dat:m2:distressed-sovereign-and-bank-capital-credit:at1}
Le jugement du Tribunal administratif fédéral du 1\up{er} octobre 2025 a annulé l'ordonnance de la FINMA
du 19 mars 2023 réduisant à zéro environ 16.5~milliards de CHF d'obligations AT1~; la FINMA a annoncé le
15 octobre 2025 qu'elle ferait appel devant le Tribunal fédéral. Aucune décision définitive n'avait été
rendue au moment de la rédaction de ce chapitre.
\end{dated}

\begin{example}[Deux ordres de pertes]\label{ex:m2:distressed-sovereign-and-bank-capital-credit:stack}
Une banque fictive dispose de 45 de fonds propres ordinaires, 16 d'AT1, 12 de rang 2 et 80 de dette senior
éligible au \omterm{def:m2:distressed-sovereign-and-bank-capital-credit:pon}{renflouement interne} (bail-in), en milliards. Une perte de 20 absorbée dans l'ordre
de la déclaration de l'UE laisse 25 aux fonds propres et toutes les obligations intactes. Une dépréciation de l'AT1 selon un
déclencheur contractuel supprime les 16 de l'AT1 et laisse les fonds propres à 45.
\end{example}

\section{Tutoriel~: une offre d'échange face au refus de participer}\label{tut:m2:distressed-sovereign-and-bank-capital-credit:recovery}

\begin{tutorial}
\textbf{Objectif.} Évaluer une \omterm{def:m2:distressed-sovereign-and-bank-capital-credit:distressed}{offre d'échange} aux rendements de sortie, évaluer le refus de participer à l'aide d'un
arbre de probabilités, trouver la participation à partir de laquelle le refus commence à rapporter, et voir
ce que change une \omterm{def:m2:distressed-sovereign-and-bank-capital-credit:cac}{clause d'action collective} agrégée.
\textbf{État final~:} \cref{fig:m2:distressed-sovereign-and-bank-capital-credit:offer,fig:m2:distressed-sovereign-and-bank-capital-credit:holdout},
\cref{ex:m2:distressed-sovereign-and-bank-capital-credit:offer,ex:m2:distressed-sovereign-and-bank-capital-credit:band,ex:m2:distressed-sovereign-and-bank-capital-credit:stack}
et les chiffres du problème du week-end.
\begin{tutsteps}
\item \textbf{L'offre}~: valeurs des obligations à un rendement de sortie, décotes, le paquet.
\omcode{code/firm/recovery/firm_recovery.py}{13}{30}{Valeurs des obligations, décotes et paquet d'échange.}\label{lst:m2:distressed-sovereign-and-bank-capital-credit:offer}
\item \textbf{Le refus de participer}~: l'arbre de probabilités, la clause, la participation d'équilibre
et la pile d'absorption des pertes.
\omcode{code/firm/recovery/firm_recovery.py}{33}{69}{L'arbre du refus de participer, la clause d'action collective, l'équilibre et l'absorption des pertes.}\label{lst:m2:distressed-sovereign-and-bank-capital-credit:holdout}
\item \textbf{Exécuter} \texttt{restructuring\_demo.tutorial()},
\texttt{restructuring\_demo.capital()} et \texttt{fig\_restructuring.py}.
\end{tutsteps}
\textbf{Pour aller plus loin.} Remplacer la clause agrégée par des clauses série par série
sur quatre obligations, et laisser un \omterm{def:m2:distressed-sovereign-and-bank-capital-credit:cac}{créancier récalcitrant} acheter 26~\% d'une série~: trouver la participation
dont le souverain a désormais besoin. Ensuite, faire dépendre la probabilité d'être payé rapidement
de la trésorerie du souverain plutôt que de la seule participation.
\end{tutorial}

\section{Construction~: évaluation de l'offre de restructuration}\label{bld:m2:distressed-sovereign-and-bank-capital-credit:recovery}

\begin{build}
\bfield{Rôle} Le desk de \omterm{def:m2:distressed-sovereign-and-bank-capital-credit:distressed}{dette en difficulté} de la petite entreprise décide s'il faut apporter ses titres à l'offre
ou refuser de participer, et son desk de crédit évalue les obligations de capital bancaire selon les deux ordres de pertes.
\bfield{Interface} \texttt{bond\_pv(face, coupon, years, y)}~; \texttt{npv\_haircut(value,
claim)}~; \texttt{Offer(new\_face, coupon, years, cash).value(y)}~;
\texttt{Holdout(paid\_soon, lawsuit\_prob, lawsuit\_pv, stuck, power).value(p)}~;
\texttt{holdout\_payoff(p, holdout, bound\_value, cac)}~;
\texttt{breakeven\_participation(holdout, tender\_value)}~; \texttt{absorb(stack, loss)}.
\bfield{Règles} Coupons annuels~; valeurs pour 100 de nominal ancien à un seul rendement de sortie~; la chance
d'être payé rapidement est la participation élevée à une puissance~; une CAC lie tous les détenteurs
au niveau de son seuil ou au-dessus, sans le versement en numéraire de l'offre~; les pertes sont absorbées dans l'ordre de la pile.
\bfield{Tests d'acceptation} \texttt{code/firm/recovery/tests/}~: une obligation dont le rendement est son taux de coupon
vaut au pair~; la valeur de l'offre baisse lorsque le rendement de sortie augmente~; la valeur du refus de participer croît
avec la participation et la clause s'applique~; la participation d'équilibre résout $H(p) = T$~; la pile
absorbe dans l'ordre.
\bfield{Pour aller plus loin} Clauses série par série et minorités de blocage~; un jeu entre créanciers
avec une distribution de croyances sur la participation~; analyse du recouvrement de la structure du capital d'une entreprise
dans un plan de faillite~; évaluation de l'AT1 avec annulation du coupon et
risque de prolongation (livre~6).
\end{build}

\begin{omsources}
\item FINMA, communiqués de presse des 19 et 23 mars 2023 et du 15 octobre 2025~; Tribunal administratif
fédéral suisse, communiqué de presse du 14 octobre 2025.
\item UBS, communiqué de presse du 19 mars 2023.
\item SRB, EBA et ECB Banking Supervision, déclaration du 20 mars 2023.
\item Comité de Bâle sur le contrôle bancaire, communiqué de presse du 13 janvier 2011 sur la capacité
d'absorption des pertes au point de non-viabilité.
\item BIS Quarterly Review, décembre 2012, encadré sur le droit applicable et la restructuration
de la dette grecque.
\item FMI, «~Strengthening the contractual framework to address collective action
problems in sovereign debt restructuring~», octobre 2014.
\end{omsources}

\section{Exercices}

\begin{exercise}[$\star$]\label{exo:m2:distressed-sovereign-and-bank-capital-credit:1}
Pour 100 d'anciennes obligations, une offre donne 60 de nouvelles obligations à 10 ans à 5~\%. Calculer sa
décote en valeur nominale et sa décote en valeur actuelle à un rendement de sortie de 8~\%~?
\end{exercise}

\begin{exercise}[$\star$]\label{exo:m2:distressed-sovereign-and-bank-capital-credit:2}
Quelle est la différence entre une \omterm{def:m2:distressed-sovereign-and-bank-capital-credit:cac}{clause d'action collective} à séries distinctes et une
clause agrégée à limbe unique~?
\end{exercise}

\begin{exercise}[$\star$]\label{exo:m2:distressed-sovereign-and-bank-capital-credit:3}
Nommer les deux déclencheurs qu'une obligation AT1 comporte, et ce dont dépend chacun d'eux.
\end{exercise}

\begin{exercise}[$\star\star$]\label{exo:m2:distressed-sovereign-and-bank-capital-credit:4}
Avec le seuil grec des deux tiers au lieu de 75~\%, sur quelle plage de participation
le refus de l'offre (holding out) est-il rentable dans l'exemple du chapitre~?
\end{exercise}

\begin{exercise}[$\star\star$]\label{exo:m2:distressed-sovereign-and-bank-capital-credit:5}
Pourquoi la lecture de la \omterm{def:m2:distressed-sovereign-and-bank-capital-credit:pp}{clause pari passu} fondée sur le paiement au prorata a-t-elle
rendu plus difficile la restructuration argentine, et comment les contrats obligataires récents y ont-ils répondu~?
\end{exercise}

\begin{exercise}[$\star\star$]\label{exo:m2:distressed-sovereign-and-bank-capital-credit:6}
Pourquoi la Grèce a-t-elle remboursé intégralement en 2012 ses \omterm{def:m2:distressed-sovereign-and-bank-capital-credit:cac}{créanciers récalcitrants} soumis au
droit anglais, alors que les détenteurs d'obligations soumises au droit grec ont subi une lourde perte~?
\end{exercise}

\begin{exercise}[$\star\star\star$]\label{exo:m2:distressed-sovereign-and-bank-capital-credit:7}
\emph{Codage.} Porter la probabilité de succès de l'action en justice de 0.3 à 0.5 et
déterminer la nouvelle participation de seuil de rentabilité. Qu'est-ce qu'une position juridique
plus forte change pour la bande des passagers clandestins~?
\end{exercise}

\begin{exercise}[$\star\star\star$]\label{exo:m2:distressed-sovereign-and-bank-capital-credit:8}
\emph{Repérer l'erreur.} «~Les obligations AT1 sont classées au-dessus des actions, donc les détenteurs d'AT1 ne peuvent pas tout perdre
alors que les actionnaires conservent quelque chose.~» Corriger cette affirmation.
\end{exercise}

\section{Problème~: Refuser ou céder}

\begin{problem}[{Problème du week-end --- la bande des passagers clandestins}]\label{pb:m2:distressed-sovereign-and-bank-capital-credit:1}
Un État en défaut offre, pour 100 d'anciennes obligations, 50 de nouvelles obligations à 15 ans à 4~\% et 5
en espèces pour ceux qui apportent leurs titres. Les nouvelles obligations devraient se traiter à un rendement de 9~\%. Un fonds
détenant 100 millions d'anciennes obligations envisage de refuser l'offre~: l'État pourrait rembourser
intégralement les \omterm{def:m2:distressed-sovereign-and-bank-capital-credit:cac}{créanciers récalcitrants} l'année suivante, avec une probabilité égale à la participation élevée à la
puissance huit~; sinon, une action en justice a 30~\% de chances de recouvrer 130 pour 100 dans six
ans~; dans le cas contraire, les obligations restent en défaut et valent 15. Les obligations comportent
une \omterm{def:m2:distressed-sovereign-and-bank-capital-credit:cac}{clause d'action collective} agrégée à limbe unique avec un seuil de 75~\%.

\medskip\noindent\textbf{Partie I --- L'offre.}
\begin{enumerate}
\item Quelle est la valeur de l'offre pour 100, avec et sans les espèces~?
\item Quelles sont ses décotes en valeur nominale et en valeur actuelle~?
\item Quelle est sa valeur à des rendements de sortie de 6~\% et de 12~\%~?
\item Pourquoi le rendement de sortie compte-t-il autant~?
\item Que reçoit le fonds s'il apporte ses 100 millions d'obligations~?
\end{enumerate}

\medskip\noindent\textbf{Partie II --- Refuser l'offre.}
\begin{enumerate}[resume]
\item Évaluer chaque branche de l'arbre.
\item Que vaut le refus de l'offre si personne n'apporte ses titres, et à des participations de 60~\% et de 74~\%~?
\item À partir de quelle participation le refus commence-t-il à être rentable~?
\item Que se passe-t-il pour un \omterm{def:m2:distressed-sovereign-and-bank-capital-credit:cac}{créancier récalcitrant} à une participation de 75~\% ou plus~?
\item Que vaudrait le refus de l'offre à une participation de 90~\% sans clause~?
\end{enumerate}

\medskip\noindent\textbf{Partie III --- Stratégie.}
\begin{enumerate}[resume]
\item Si chaque créancier raisonnait comme le fonds, la participation pourrait-elle se stabiliser entre 60.9~\%
et 75~\%~?
\item Pourquoi les espèces versées pour l'apport des titres comptent-elles une fois la clause contraignante~?
\item Comment le fonds pourrait-il bloquer la clause, et pourquoi est-ce plus difficile avec une agrégation~?
\item Qu'ont reçu les \omterm{def:m2:distressed-sovereign-and-bank-capital-credit:cac}{créanciers récalcitrants} grecs soumis au droit anglais, et pourquoi~?
\item Qu'ont subi les détenteurs argentins ayant accepté l'\omterm{def:m2:distressed-sovereign-and-bank-capital-credit:distressed}{offre d'échange} à cause de la victoire des \omterm{def:m2:distressed-sovereign-and-bank-capital-credit:cac}{créanciers récalcitrants}~?
\end{enumerate}

\medskip\noindent\textbf{Partie IV --- Jugement.}
\begin{enumerate}[resume]
\item Pourquoi la probabilité d'être payé rapidement est-elle difficile à estimer~?
\item Un fonds qui achète à bas prix des obligations en défaut pour engager des poursuites doit-il pouvoir bloquer
une restructuration~?
\item Comment couvrir la position du fonds pendant qu'il attend~?
\item Énoncer le \emph{résultat nommé}~: la bande de participation
dans laquelle le refus est rentable, et ce que la clause fait au-dessus.
\item En une phrase~: quel problème les \omterm{def:m2:distressed-sovereign-and-bank-capital-credit:cac}{clauses d'action collective} résolvent-elles~?
\end{enumerate}
\end{problem}

\section{Questions d'entretien}

\begin{interviewq}[$\star$ \iqroles{trader, chercheur}]\label{iq:m2:distressed-sovereign-and-bank-capital-credit:1}
En quoi une restructuration de dette souveraine diffère-t-elle d'une faillite d'entreprise~?
\end{interviewq}

\begin{interviewq}[$\star$ \iqroles{trader, banque}]\label{iq:m2:distressed-sovereign-and-bank-capital-credit:2}
Qu'est-ce qu'une obligation AT1, et quels sont ses risques~?
\end{interviewq}

\begin{interviewq}[$\star\star$ \iqroles{chercheur}]\label{iq:m2:distressed-sovereign-and-bank-capital-credit:3}
Comment évalueriez-vous une \omterm{def:m2:distressed-sovereign-and-bank-capital-credit:distressed}{offre d'échange} en situation de détresse~?
\end{interviewq}

\begin{interviewq}[$\star\star$ \iqroles{trader, chercheur}]\label{iq:m2:distressed-sovereign-and-bank-capital-credit:4}
Expliquez le problème des \omterm{def:m2:distressed-sovereign-and-bank-capital-credit:cac}{créanciers récalcitrants} et la manière dont les \omterm{def:m2:distressed-sovereign-and-bank-capital-credit:cac}{clauses d'action collective} y répondent.
\end{interviewq}

\begin{interviewq}[$\star\star$ \iqroles{risque, banque}]\label{iq:m2:distressed-sovereign-and-bank-capital-credit:5}
Qu'a appris aux investisseurs la dépréciation totale des AT1 de Credit Suisse sur les fonds propres bancaires~?
\end{interviewq}

\begin{interviewq}[$\star\star\star$ \iqroles{développeur}]\label{iq:m2:distressed-sovereign-and-bank-capital-credit:6}
Concevez un outil qui suit, pour les obligations en circulation d'un État souverain, le droit applicable, le
type de \omterm{def:m2:distressed-sovereign-and-bank-capital-credit:cac}{clause d'action collective} et le nombre de voix nécessaires pour engager chaque série.
\end{interviewq}
